% ---------------------------------------------------------------- 01-introduction.tex
\section{Introduction}\label{sec:intro}

Our starting point is the two-state firm of Lecture~1 (pp.~9, 47--48):\footnote{``Lecture'' refers to the MH6806 course slides \citep{vaidya2026}; page numbers are those of the PDF files.} financed by equity alone or by debt and equity, it pays the same 120 or 60, so buying the cheaper package and shorting the dearer one would be a money machine whenever their prices differed. \citet{mm1958} turn this argument into the benchmark result: when capital markets are perfect and the firm's investment policy is held fixed, financing decisions cannot change firm value:
\begin{equation}
\VL = \VU. \label{eq:mm}
\end{equation}

Debt and equity merely partition the same operating cash flows; to borrow a familiar image, the size of the pizza does not depend on how it is sliced. The proposition is valuable less as a description of the world than as a benchmark: if capital structure affects value in practice, at least one of its assumptions must fail, and the useful question becomes which one, and by how much.

We study the assumption with the most direct consequences for practice: the absence of corporate taxation. Because interest is deductible in almost every corporate tax system, debt reduces the government's claim on operating income, and \citet{mm1963} show that this alone makes financing matter. Corporate taxation is also the friction that policy changes most directly. The 2017 US Tax Cuts and Jobs Act lowered the federal rate from 35\% to 21\% and capped interest deductions, while Singapore combines a 17\% corporate rate with broad exemptions for investment income at the investor level. Lecture~1 (p.~50) lists tax shields among the benefits of debt and financial distress among its costs; we put numbers on that trade-off. Our research question is therefore: \textit{does the corporate income tax turn debt into a free lunch?}

We judge every result against two conditions. (i)~\textit{Attribution}: any value created by debt must vanish as the tax rate goes to zero, so that it can be attributed to the tax. Inside our model this holds by construction, because tax is the only benefit of debt; outside it, S-REITs, which pay almost no tax but still borrow about 40\% of their assets, show that leverage alone is not evidence of a tax benefit (Section~\ref{sec:sg}). (ii)~\textit{Robustness}: the shape of the result, a gain up to a debt limit followed by a loss, must survive reasonable changes in parameters; findings that depend on a particular parameter value are reported as properties of the model rather than of the world.

Our analysis draws on a well-established literature. \citet{mm1963} show that with corporate taxes $\VL = \VU + \tauc D$, which taken literally implies all-debt financing. The trade-off theory bounds this result with bankruptcy costs \citep{kraus1973}. \citet{miller1977} shows that personal taxes on interest income can offset the corporate advantage, and \citet{deangelo1980} show that the marginal value of interest deductions falls when they cannot always be used. Empirically, \citet{graham2000} estimates the capitalised tax benefit of debt at 9.7\% of firm value and documents that many firms borrow conservatively, while distress costs are estimated at 10--20\% of firm value once distress occurs \citep{andrade1998} and, in risk-adjusted present value, at about 4.5\% for a BBB-rated firm \citep{almeida2007}.

Our answer is no: the corporate tax makes debt valuable to shareholders, but the gain is paid for by the government and capped by the firm's weakest year, so it is not free. We follow this mechanism from the two-state firm to calibrated settings. Sections~\ref{sec:model}--\ref{sec:robust} present the base model and test its robustness. Section~\ref{sec:base} also reruns the law-of-one-price trade with taxes and reads the result with CAPM betas, WACC and a stress test (Lectures~1, 2 and~5). Section~\ref{sec:ext} relaxes, one at a time, the four simplifications most likely to overturn the result. Section~\ref{sec:apps} confronts the model with the 2017 US tax reform, the Singapore tax system and the reported cash flows of two Singapore-listed firms. Section~\ref{sec:concl} concludes, and Table~\ref{tab:summary} summarises which results survive.
